% Example: Fractions, square roots, and exponents
% Demonstrates \frac for fractions, \sqrt for roots (including nth roots),
% and superscript/subscript notation for exponents and indices.
\documentclass[12pt, a4paper]{article}
\usepackage{amsmath, amssymb, amsthm}
\usepackage{fontspec}
\begin{document}

\section*{Fractions, Roots, and Exponents}

\subsection*{Fractions}

A simple fraction:

\[
  \frac{1}{2} + \frac{1}{3} = \frac{5}{6}
\]

Nested fractions:

\[
  \frac{1}{1 + \frac{1}{1 + \frac{1}{2}}}
\]

\subsection*{Square and nth Roots}

Square root:

\[
  \sqrt{2} \approx 1.4142
\]

nth root (using the optional argument):

\[
  \sqrt[3]{8} = 2, \qquad \sqrt[n]{x^n} = |x|
\]

Root of a fraction:

\[
  \sqrt{\frac{a^2 + b^2}{2}}
\]

\subsection*{Exponents and Subscripts}

\[
  a^2 + b^2 = c^2
\]

\[
  2^{10} = 1024, \qquad e^{i\pi} = -1
\]

\[
  x_1 + x_2 + \cdots + x_n = \sum_{i=1}^{n} x_i
\]

Combined fraction, root, and exponent:

\[
  \left( \frac{x^2 + 1}{\sqrt{x}} \right)^{3}
\]

\end{document}
